\documentclass[12pt,a4paper]{article}
\usepackage{cmap} % Makes the PDF copiable. See http://tex.stackexchange.com/a/64198/25761
\usepackage[T1]{fontenc}
\usepackage[brazil]{babel}
\usepackage[utf8]{inputenc}
\usepackage{amsmath}
\usepackage{amsfonts}
\usepackage{amssymb}
\usepackage{amsthm}
\usepackage[usenames,svgnames,dvipsnames]{xcolor}
\usepackage{hyperref}
\usepackage{multicol}
\usepackage{graphicx}
\usepackage{cancel}
\usepackage[top=2cm, bottom=2cm, left=2cm, right=2cm]{geometry}

\hypersetup{
    colorlinks = true,
    allcolors = {blue}
}

% TODO: Consider using exsheets
% http://linorg.usp.br/CTAN/macros/latex/contrib/exsheets/exsheets_en.pdf
%
% http://ctan.org/tex-archive/macros/latex/contrib/exercise/
% Options: answerdelayed,lastexercise,noanswer
\usepackage[answerdelayed,lastexercise]{exercise}

\addto\captionsbrazil{%
\def\listexercisename{Lista de exerc\'icios}%
\def\ExerciseName{Exerc\'icio}%
\def\AnswerName{Solu\c{c}\~ao do exerc\'icio}%
\def\ExerciseListName{Ex.}%
\def\AnswerListName{Solu\c{c}\~ao}%
\def\ExePartName{Parte}%
\def\ArticleOf{de\ }%
}

\renewcommand{\ExerciseHeaderTitle}{(\ExerciseTitle)\ }
\renewcommand{\ExerciseListHeader}{%\ExerciseHeaderDifficulty%
\textbf{%\ExerciseListName\
\ExerciseHeaderNB.\ %
%\ --- \
\ExerciseHeaderTitle}%
%\ExerciseHeaderOrigin
\ignorespaces}
\renewcommand{\AnswerListHeader}{\textbf{\ExerciseHeaderNB.\ (\AnswerListName)\ }}

\newcommand*\diff{\mathop{}\!\mathrm{d}}
\newcommand*\sen{\operatorname{sen}}
\newcommand*\R{\mathbb{R}}

\renewcommand{\theenumi}{\alph{enumi}}
\renewcommand\labelenumi{(\theenumi) }

\newcommand*\tipo{Prova III}
\newcommand*\turma{NEXM241-D}
\newcommand*\disciplina{CDI1001}
\newcommand*\eu{Helder G. G. de Lima}
\newcommand*\data{08/06/2024}

\author{\eu}
\title{\tipo - \disciplina}
\date{\data}

\begin{document}
\thispagestyle{empty}
\newgeometry{margin=2cm,bottom=0.5cm}
\begin{center}
\includegraphics[width=9.0cm]{marca} \\
\textbf{\tipo\ (\disciplina / \turma)} \\
Prof. \eu\footnote{
Este é um material de acesso livre distribuído sob os termos da licença \href{https://creativecommons.org/licenses/by-sa/4.0/deed.pt_BR}{Creative Commons BY-SA 4.0}}
\end{center}

\noindent Nome do(a) aluno(a): \underline{\hspace{9,7cm}} Data: \underline{\data}

%\section*{Instruções}
\begin{center}\fbox{
\begin{minipage}{14cm}

{\footnotesize
\begin{itemize}
\renewcommand{\theenumi}{\Roman{enumi}}
\item Identifique-se em todas as folhas.
\item Mantenha o celular e os demais equipamentos eletrônicos desligados durante a prova.
\item Justifique cada resposta com cálculos ou argumentos baseados na teoria estudada.
\item Resolva $5$ das $6$ questões (deixe claro que questão não deverá ser corrigida).
\end{itemize}
}

\end{minipage}
}
\end{center}

%\section*{Questões}
\begin{ExerciseList}
\Exercise[title={2,0}] Seja $y = f(x) = e^{2x} \cos(x)$. Verifique que $y^{\prime\prime} = 4 y^\prime - 5 y$.
\Answer Pela regra do produto e pela regra da cadeia, a derivada de $f$ é dada por:
\[
    y^\prime
    = (e^{2x} \cos(x))^\prime
    = (e^{2x})^\prime \cos(x) + e^{2x}(\cos(x))^\prime
    = 2e^{2x} \cos(x) - e^{2x} \sen(x).
\]
Além disso, a segunda derivada de $f$ é dada por:
\begin{align*}
    y^{\prime\prime}
    & = (2e^{2x} \cos(x) - e^{2x} \sen(x))^\prime \\
    & = 2(e^{2x} \cos(x))^\prime - (e^{2x} \sen(x))^\prime \\
    & = 2[2e^{2x} \cos(x) - e^{2x} \sen(x)]
      - [(e^{2x})^\prime \sen(x) + e^{2x}(\sen(x))^\prime] \\
    & = 4e^{2x} \cos(x) - 2e^{2x} \sen(x)
    - 2e^{2x} \sen(x) - e^{2x}\cos(x) \\
    & = 3e^{2x} \cos(x) - 4e^{2x} \sen(x).
\end{align*}
Portanto,
\begin{align*}
    4 y^\prime - 5 y
    & = 4[2e^{2x} \cos(x) - e^{2x} \sen(x)] - 5[e^{2x} \cos(x)] \\
    & = 8e^{2x} \cos(x) - 4e^{2x} \sen(x) - 5e^{2x} \cos(x) \\
    & = 3e^{2x} \cos(x) - 4e^{2x} \sen(x) = y^{\prime\prime}.
\end{align*}

\Exercise[title={2,0}] Utilize uma aproximação linear para estimar o valor de $\sqrt[3]{1003}$.
\Answer Considerando $g(x) = \sqrt[3]{x} = x^{\frac{1}{3}}$, tem-se $\sqrt[3]{1003} = g(1003)$. Além disso, como $g$ é uma função derivável em todo $x \neq 0$ e $g^\prime(x) = \frac{1}{3}x^{-\frac{2}{3}} = \dfrac{1}{3\sqrt[3]{x^2}}$, resulta que para valores de $x$ próximos de $x_0 = 1000 = 10^3$, tem-se a aproximação linear
\[
    g(x) \approx g(1000) + g^\prime(1000)(x - 1000).
\]
Então, como $g(1000) = \sqrt[3]{1000} = 10$ e $g^\prime(1000) = \dfrac{1}{3\sqrt[3]{1000^2}} = \dfrac{1}{3\sqrt[3]{(10^3)^2}} = \dfrac{1}{3\cdot 10^2} = \dfrac{1}{300}$, conclui-se que:
\[
    \sqrt[3]{1003}
    \approx 10 + \dfrac{1}{300}(1003 - 1000)
    = 10 + \dfrac{1}{300}\cdot 3
    = 10 + \dfrac{1}{100}
    = \dfrac{1000 + 1}{100}
    = 10,01.
\]
Apenas para comparação, o valor obtido em uma calculadora seria de $\sqrt[3]{1003} \approx 10,00999$.

\Exercise[title={2,0}]
\begin{minipage}{0.6\textwidth}
    Um tanque invertido em forma de cone possui uma altura de $4m$ e um raio na base de $2m$. Ele está sendo preenchido com água a uma taxa de $6m^3/min$. Qual é a velocidade de elevação do nível da água quando sua distância até o fundo atinge $3m$?
\end{minipage}\hfill
\begin{minipage}{0.3\textwidth}
    \begin{center}
        \includegraphics[width=5cm]{img/prova-3-nex-cone.png}
    \end{center}
\end{minipage}
\Answer Para cada instante $t$ (em minutos), considere que:
\begin{itemize}
    \item A altura da água no tanque (em metros) é $h = h(t)$
    \item O raio da superfície da água (em metros) é $r = r(t)$
    \item O volume de água no tanque (em metros cúbicos) é $V = V(t)$
\end{itemize}
Então a taxa de preenchimento do tanque é a taxa de variação do volume em relação ao tempo, ou seja, $\frac{\diff{V}}{\diff{t}} = 6$, e o objetivo é calcular a taxa de variação da altura em relação ao tempo, isto é, $\frac{\diff{h}}{\diff{t}}$, quando a altura é $h=3$.

Sabe-se que o volume $V$ de um cone circular com altura $h$ e raio da base $r$ é dado por $V = \frac{1}{3}\pi r^2 h$. Além disso, considerando as dimensões do tanque ($2 m$ de raio e $4 m$ de altura), o raio $r$ da superfície da água é sempre metade da altura $h$ correspondente, isto é, $r = h/2$. Assim, o volume de água quando sua altura é $h$ é dado por $V = \frac{1}{3}\pi {\left(\frac{h}{2}\right)}^2 h = \frac{\pi}{12} h^3.$ Derivando ambos os membros dessa equação em relação a $t$, com o auxílio da regra da cadeia, obtém-se uma relação entre a taxa de variação do volume em relação ao tempo e a taxa de variação procurada:
\[
    \frac{\diff{V}}{\diff{t}}
    = \frac{\diff}{\diff{t}}\left(\frac{\pi}{12} h^3\right)
    = \frac{\pi}{12} \frac{\diff}{\diff{t}}\left(h^3\right)
    = \frac{\pi}{12} 3h^2 \frac{\diff{h}}{\diff{t}}
    = \frac{\pi}{4} h^2 \frac{\diff{h}}{\diff{t}}.
\]
Substituindo $h=3$ e $\frac{\diff{V}}{\diff{t}} = 6$, obtém-se:
\[
    \frac{\diff{V}}{\diff{t}}
    = \frac{\pi}{4} h^2 \frac{\diff{h}}{\diff{t}}
    \Rightarrow
    6
    = \frac{\pi}{4} 3^2 \frac{\diff{h}}{\diff{t}}
    \Rightarrow
    24
    = 9\pi \frac{\diff{h}}{\diff{t}}
    \Rightarrow
    \frac{\diff{h}}{\diff{t}}
    = \frac{8}{3\pi}.
\]

Portanto, o nível da água sobe a uma taxa de $\dfrac{8}{3\pi} m^3 / min$, ou pouco menos de $1 m^3 / min$ (já que $\pi$ é um pouco maior do que $3$).
\Exercise[title={2,0}] Determine $\displaystyle\lim_{x \to 0} \frac{x^2 + x - \sen(x)}{x^3 + 21x^2}$ utilizando, se possível, as regras de L'Hôpital.
\Answer Por L'Hôpital,
\[
\lim_{x \to 0} \frac{ \cancelto{0}{x^2 + x - \sen(x)} }{ \cancelto{0}{x^3 + 21x^2}}
\stackrel{\text{L.H.}}{=}\lim_{x \to 0} \frac{ \cancelto{0}{2x + 1 - \cos(x)} }{ \cancelto{0}{3x^2+42x} }
\stackrel{\text{L.H.}}{=}\lim_{x \to 0} \frac{2 + \sen(x)}{ 6x + 42}
=\frac{2 + 0}{ 0 + 42}
=\frac{1}{21}.
\]

\Exercise[title={2,0}] Calcule $\displaystyle\lim_{x\to 0^+} \left[\sen(x)\right]^{\left[\frac{1}{\ln(x)}\right]}$ usando, se possível, as regras de L'Hôpital.
\Answer Seja $L = \lim_{x\to 0^+} \left[\sen(x)\right]^{\left[\frac{1}{\ln(x)}\right]}$. Então, aplicando o logaritmo em ambos os membros,
\begin{align*}
    \ln(L)
    & = \ln\left(\lim_{x\to 0^+} \left[\sen(x)\right]^{\left[\frac{1}{\ln(x)}\right]}\right)
    = \lim_{x\to 0^+} \ln\left(\left[\sen(x)\right]^{\left[\frac{1}{\ln(x)}\right]}\right)\\
    & = \lim_{x\to 0^+} \left[\frac{1}{\ln(x)}\right] \ln\left(\sen(x)\right)
    = \lim_{x\to 0^+} \frac{\ln\left(\sen(x)\right)}{\ln(x)}.
\end{align*}
Por L'Hôpital,
\begin{align*}
    \ln(L)
    &
    = \lim_{x\to 0^+} \frac{(\ln\left(\sen(x)\right))^\prime}{(\ln(x))^\prime}
    = \lim_{x\to 0^+} \frac{\left(\frac{\cos(x)}{\sen(x)}\right)}{\left(\frac{1}{x}\right)}
    = \lim_{x\to 0^+} \frac{\cos(x)}{\sen(x)}\cdot x\\
    &
    = \lim_{x\to 0^+} \cos(x)\frac{x}{\sen(x)}
    = \lim_{x\to 0^+} \frac{\cos(x)}{\frac{\sen(x)}{x}}
    = \frac{1}{1}
    = 1
\end{align*}
Portanto, como $\ln(L) = 1$, conclui-se que $L = e^1 = e$.

\Exercise[title={2,0}] Para cada item abaixo, diga se é verdadeiro (V) ou falso (F).

\emph{Observação 1:} Não é necessário justificar.

\emph{Observação 2:} Cada item marcado errado anula a pontuação de um item marcado corretamente. Você tem a opção de deixar o item em branco e, nesse caso, você nem ganha e nem perde pontos.

\begin{enumerate}
  \item {\bf ( \ \ )} \ Se $p(x) = x^5 + x^2 - x$ então existe algum ponto $c \in (-1, 1)$ tal que $p^\prime(c) = 0$.
  \item {\bf ( \ \ )} \ Se $g(x) = 3^x - x^2$ então existe algum ponto $c \in (0, 2)$ tal que $g^\prime(c) = 2$.
  \item {\bf ( \ \ )} \ A função $f(x) = \cos(x)$ não tem ponto de máximo global;
  \item {\bf ( \ \ )} \ A função $q(x) = 2x + \sen(x)$ é crescente em todos os intervalos $(a, b) \subset \R$;
\end{enumerate}

\Answer \begin{enumerate}
  \item \textbf{Verdadeiro}: Como $p(x) = x^5 + x^2 - x$ é contínua em $[-1, 1]$ e derivável em $(-1, 1)$, e $p(-1) = 1 = p(1)$, o teorema de Rolle garante que existe algum ponto $c \in (-1, 1)$ tal que $p^\prime(c) = 0$.
  \item \textbf{Verdadeiro}: Como $g(x) = 3^x - x^2$ é contínua em $[0, 2]$ e derivável em $(0, 2)$, o teorema do valor médio de Lagrange garante que existe algum ponto $c \in (0, 2)$ tal que
  \[
  g^\prime(c)
  = \frac{g(2) - g(0)}{2 - 0}
  = \frac{(9-4) - (1-0)}{2}
  = \frac{5 - 1}{2}
  = \frac{4}{2}
  = 2.
  \]
  \item \textbf{Falso}: Todo ponto da forma $x = 2 \pi n$, com $n \in \mathbb{Z}$ é um ponto de máximo global de $f(x) = \cos(x)$, pois $\cos(2 \pi n) = 1$ e não ocorre $\cos(x) > 1$ em nenhum outro $x \in \R$.
  \item \textbf{Verdadeiro}: Como $q(x) = 2x + \sen(x)$ tem derivada $q^\prime(x) = 2 + \cos(x)$ e $\cos(x) \geq -1$, para todo $x \in \R$, tem-se então $q^\prime(x) \geq 1 > 0$ e $q(x)$ é crescente em todo intervalo $(a, b) \subset \R$.
\end{enumerate}
\end{ExerciseList}

\begin{center}
BOA PROVA!
\end{center}

\newpage
\restoregeometry
\section*{Respostas}
\shipoutAnswer
\end{document}
